% !TEX root = ../arbitrage.tex
\section{Mechanism Design: Equilibrium Redesign}
\label{sec:mechanism}

The preceding sections identify a precise failure mode. Under cheap-talk costs,
anthropomorphic marketing can pool across governance types, users rationally
invest under the prior, and regulators abstain when inspection is too costly.
The policy implication is not that regulators should decide whether AI systems
are persons. It is that the public signal environment should be redesigned so
that inconsistent ontological claims no longer remain costless. In
mechanism-design terms, the problem is to make truthful revelation of $\theta_F$ incentive
compatible rather than to ask firms to volunteer it \citep{myerson1979}.

The proposals below therefore alter the game in Section~\ref{sec:bayesian} at
three different points: the cost of inconsistent messages, the cost of
anthropomorphic claims, and the regulator's cost of verification. Each proposal
has the same structure. It modifies a model parameter, states the equilibrium
shift induced by that modification, and points to an existing legal form that
establishes feasibility. The analogs are not claims that current doctrine already
solves AI ontology. They show that adjacent legal systems already know how to
attach consequences to half-truths, require substantiation before marketing, and
make internal records available to oversight institutions.

\subsection{Sanctionable Inconsistency}
\label{subsec:sanctionable-inconsistency}

The first intervention treats the firm's public message as a joint object rather
than as two isolated texts. In the cheap-talk model, marketing can send $a$
while policy and liability documents send $p$, with no cost assigned to the
inconsistency. A sanctionable-inconsistency rule would require firms to maintain
a claim register linking outward-facing marketing, safety statements, terms of
service, incident disclosures, and liability positions. If the firm represents
the system as agentic, caring, autonomous, or safety-critical in one channel
while denying the legal significance of the same representation in another,
that cross-channel inconsistency becomes independently sanctionable.

\paragraph{Parameter changed.}
This adds a penalty term $S_I$ to the firm's payoff for inconsistent message
pairs. If $\rho_R$ is the probability that $R$ detects or accepts proof of the
inconsistency, the cheap-talk premium becomes
\[
\Delta_F-\rho_R S_I .
\]
Pooling on anthropomorphic marketing is no longer profitable when
$\rho_R S_I>\Delta_F$. The rule therefore attacks the exact term that made
arbitrage attractive: the rent earned from pricing the system as agent in the
user channel while pricing it as tool in the liability channel.

\paragraph{Equilibrium shift.}
The induced shift is not necessarily full separation. It can also operate by
constraining off-path beliefs. If $R$ observes an inconsistent pair $(a,p)$,
then the admissible posterior places greater mass on $\theta_F=\mathsf{L}$ or
on high opacity $\theta_S$. Firms anticipating that posterior either harmonize
their claims or choose the deflationary message across channels. In both cases,
the cheap-talk pooling equilibrium loses its most valuable deviation: firms can
no longer take the marketing upside of $a$ while preserving the liability
downside of $p$.

\paragraph{Existing-law analog.}
The legal analog is familiar. FTC Act Section 5 treats unfair or deceptive acts
or practices as sanctionable, and advertising-substantiation doctrine already
asks whether the net impression of a claim is misleading to the relevant
audience. Securities antifraud rules provide the parallel logic for investor
communications: Rule 10b-5 targets material misstatements and omissions that
make statements misleading in context. The proposed rule imports that
cross-document logic without importing securities doctrine wholesale. Its point
is narrower: a firm should not be able to produce one ontology for users and an
incompatible ontology for adjudicators without paying a regulatory price.

\subsection{Costly Signaling Through Third-Party Audits}
\label{subsec:costly-signaling-audits}

The second intervention makes anthropomorphic or agency-adjacent marketing a
costly signal. A firm may still describe a system as a companion, agent,
co-worker, tutor, or safety-critical assistant, but only after a qualified third
party verifies the governance conditions that make such a claim non-misleading:
evaluation coverage, red-team results, escalation pathways, incident-response
capacity, model-update controls, and post-deployment monitoring. The audit does
not certify consciousness or moral status. It certifies the institutional
conditions behind the claim.

\paragraph{Parameter changed.}
This implements the audit-trail cost from Proposition~\ref{prop:cheap-talk-separation}.
Let audit intensity be $e$, with $c_H(e)=k_He$ for high-governance firms and
$c_L(e)=k_Le$ for low-governance firms, where $0<k_H<k_L$. The separating
interval derived in Section~\ref{subsec:audit-trail-separation} is
\[
\frac{G}{k_L}\leq e\leq \frac{G}{k_H}.
\]
Above the lower threshold, low-governance firms cannot profitably mimic. Below
the upper threshold, high-governance firms can still afford to signal. The
regulator's task is therefore not to maximize audit burden. It is to set $e$ in
the interval where signal cost separates types.

\paragraph{Equilibrium shift.}
The shift is from pooling to separating PBE. High-governance firms send audited
anthropomorphic claims; low-governance firms either invest in governance until
their cost curve changes or use deflationary marketing. Users can then condition
investment on an informative signal rather than on undisciplined warmth in
product copy. Regulators can condition abstention on the audit result rather
than on an unverified prior. This is the Spence logic of costly signaling
applied to AI governance claims \citep{spence1973}; it also follows the
advertising-as-signal insight that market messages become informative only when
false messages are sufficiently expensive \citep{milgrom_roberts1986}.

\paragraph{Existing-law analog.}
The analogs are pharmaceutical efficacy review and food-labeling controls. In
both domains, public-facing claims are not left entirely to seller narrative:
claims about therapeutic effect or product identity must be supported before
they can be used to shape consumer reliance. The same design can be translated
to AI without pretending that AI companionship is a drug or a food. The common
institutional pattern is substantiation-before-reliance: when a claim predictably
changes how vulnerable users act, the law may require evidence before allowing
the claim to be marketed at scale.

\subsection{Auditable Records for Regulator-Readable Verification}
\label{subsec:auditable-records}

The third intervention reduces the regulator's verification cost by requiring
append-only, regulator-readable records tied to capability and relational
claims. Records should include the claim identifier, model version, relevant
evaluation suite, deployment date, material model updates, safety interventions,
known incident classes, and the internal owner responsible for the claim. The
point is not to publish every interaction or expose user data. It is to create a
durable evidentiary substrate that lets $R$ test whether a public claim remained
true across model updates and incidents.

\paragraph{Parameter changed.}
In the pooling condition from Section~\ref{subsec:cheap-talk-pooling}, the
regulator abstains when $K_{\theta_R}>\pi(\mathsf{L})D_R$. Auditable records
lower $K_{\theta_R}$ and can increase expected detection value $D_R$:
\[
K'_{\theta_R}=K_{\theta_R}-\lambda,\qquad
D'_R=D_R+\eta .
\]
When $K'_{\theta_R}\leq \pi(\mathsf{L})D'_R$, inspection becomes sequentially
rational. Records also improve the public channel $z$ by giving journalists,
researchers, and civil society a procedural object to demand or contest, rather
than leaving them to infer governance quality from outputs alone.

\paragraph{Equilibrium shift.}
The shift is from abstention under unverifiable priors to credible inspection.
Once firms know that claims are tied to records, a low-governance firm faces a
higher expected cost from exaggeration even if it is not inspected in every
case. This complements the first two proposals. Sanctionable inconsistency tells
firms which cross-channel contradictions will be penalized; costly signaling
tells firms what evidence is required before making anthropomorphic claims;
auditable records make both rules administrable over time.

\paragraph{Existing-law analog.}
Several regulatory systems already use this form. GDPR Article 30 requires
records of processing activities. Sarbanes-Oxley Section 404 links public
reporting to internal control over financial reporting. EU AI Act Article 12
requires logging capabilities for high-risk AI systems. None is a complete
template for relational AI governance, but each shows the same institutional
move: private organizations must produce internal records in a form that public
oversight can inspect. That is the rule-of-law function of documentation in
complex technical systems \citep{hadfield2017rules}.

Taken together, these interventions change the equilibrium rather than merely
condemning the equilibrium. Sanctionable inconsistency taxes contradictory
messages. Costly audits create a separating signal for $\theta_F$. Auditable
records lower $K_{\theta_R}$ and make public signal $z$ less noisy. The result
is not metaphysical settlement. It is a governance design in which firms remain
free to choose their ontology, but not to arbitrage incompatible ontologies
across audiences at zero cost.
